\documentclass[12pt,a4paper]{article}

\usepackage[spanish]{babel}
\usepackage[utf8]{inputenc}
\usepackage{amsmath, amssymb, amsfonts}
\usepackage{graphicx}
\usepackage{geometry}
\usepackage{float}
\usepackage{booktabs}
\usepackage{hyperref}
\usepackage{setspace}
\usepackage{caption}

\geometry{margin=2.5cm}
\setstretch{1.2}
\setlength{\parskip}{1em}
\setlength{\parindent}{0pt}

%-------------------- PORTADA ---------------------
\begin{document}

\begin{titlepage}
\centering
{\large \textbf{UNIVERSIDAD DE CARABOBO}}\\[0.2cm]
{\large \textbf{FACULTAD EXPERIMENTAL DE CIENCIAS Y TECNOLOGÍA}}\\[0.2cm]
{\large \textbf{DEPARTAMENTO DE COMPUTACIÓN}}\\[1.5cm]

{\Large \textbf{Arquitectura del Computador}}\\[2cm]

{\Huge \textbf{Informe Técnico}}\\[0.8cm]
{\LARGE \textbf{Comparación de Boost Flat Map vs Google B-Tree en la Administración de Caché}}\\[2cm]

\begin{flushleft}
\textbf{Integrantes:}\\
Samuel Cordero \hspace{0.3cm} C.I. V-31.678.592\\
Dervis Martínez \hspace{0.3cm} C.I. V-31.456.xxx\\[1cm]

\textbf{Profesor:} Ing. (coloca el nombre del docente si aplica)\\[1cm]
\textbf{Sección:} (coloca la sección)\\[1cm]
\textbf{Fecha:} Octubre 2025
\end{flushleft}

\vfill

{\large Universidad de Carabobo, Venezuela}
\end{titlepage}

\tableofcontents
\newpage

%--------------------------------------------------------
\section{El Problema}

En los sistemas modernos, el rendimiento global depende en gran medida de la eficiencia con la que se gestionan los datos en memoria.  
La memoria caché desempeña un papel fundamental, pues permite acceder a información frecuente con menor latencia.

El problema planteado consiste en evaluar el rendimiento de dos estructuras de datos aplicadas a la administración de caché:  
la \textbf{Boost Flat Map} y la \textbf{Google B-Tree}.  
Cada una implementa diferentes estrategias de almacenamiento, búsqueda e inserción, lo cual puede afectar de manera significativa la velocidad y el uso de la memoria.

El objetivo principal es determinar cuál de estas estructuras ofrece mejor desempeño al simular el comportamiento de una memoria caché,  
considerando los siguientes factores:
\begin{itemize}
    \item Tiempo de búsqueda, inserción y eliminación.
    \item Latencia promedio por operación.
    \item Tasa de aciertos (\textit{hit rate}) bajo distintos patrones de acceso.
\end{itemize}

El experimento se desarrolló en C++ con código totalmente estructurado y sin el uso de librerías estándar (STL),  
implementando manualmente las estructuras requeridas.

%--------------------------------------------------------
\section{Técnicas Computacionales para su Solución}

\subsection{Diseño General}

Se implementaron dos estructuras de caché:

\begin{enumerate}
    \item \textbf{FlatMapCache:} basada en un arreglo ordenado que permite búsquedas mediante algoritmo binario.
    \item \textbf{BTreeCache:} basada en un árbol balanceado del tipo B-Tree que distribuye las claves en nodos.
\end{enumerate}

Ambas estructuras fueron desarrolladas en un único archivo fuente \texttt{cache-compare.cpp},  
definiendo además una estructura auxiliar de arreglo dinámico propia (\texttt{Vector<T>})  
para gestionar los elementos sin emplear \texttt{std::vector}.

\subsection{FlatMapCache}

El \textbf{FlatMapCache} simula el comportamiento de la estructura \texttt{Boost::flat\_map}, almacenando los pares clave–valor en un bloque contiguo de memoria.  
Esto permite:
\begin{itemize}
    \item Búsquedas rápidas mediante búsqueda binaria (\(O(\log n)\)).
    \item Inserciones más lentas, ya que requieren desplazar los elementos.
\end{itemize}

Su principal ventaja es la \textbf{alta localidad de referencia}, lo que mejora el uso de la memoria caché del procesador.

\subsection{BTreeCache}

El \textbf{BTreeCache} emula la estructura \texttt{google::btree\_map}.  
Se compone de nodos con múltiples claves que se dividen (\emph{split}) cuando se llenan,  
manteniendo el árbol balanceado y garantizando tiempos de búsqueda e inserción logarítmicos.  

Características:
\begin{itemize}
    \item Inserciones más estables que FlatMap.
    \item Búsqueda con costo \(O(\log n)\).
    \item Mayor sobrecarga por punteros (menor localidad de caché).
\end{itemize}

\subsection{Generador de Trazas (TraceGen)}

Para probar el comportamiento bajo distintos escenarios de carga se creó un módulo llamado \textbf{TraceGen},  
que genera secuencias de accesos (claves) con tres patrones posibles:

\begin{itemize}
    \item \textbf{Random:} claves elegidas de manera uniforme aleatoria.
    \item \textbf{Sequential:} claves accedidas en orden creciente y cíclico.
    \item \textbf{Zipf:} distribución sesgada donde unas pocas claves son accedidas con mayor frecuencia.
\end{itemize}

La distribución Zipf fue implementada generando una función de distribución acumulada (CDF)  
y aplicando búsqueda binaria sobre ella para obtener claves de acuerdo a su probabilidad.

\subsection{Medición de Rendimiento}

Se midieron los tiempos de ejecución en nanosegundos utilizando la librería \texttt{std::chrono}.  
Para cada estructura se evaluó:

\begin{itemize}
    \item Latencia promedio (tiempo por operación).
    \item Cantidad de aciertos en caché.
    \item Tiempo medio de inserciones y eliminaciones puras.
\end{itemize}

Los resultados se almacenaron en una estructura \texttt{Stats} con las variables:
\texttt{avg\_ns\_per\_op}, \texttt{hits} y \texttt{ops}.

%--------------------------------------------------------
\section{Análisis y Comparación de Resultados}

\subsection{Parámetros de Ejecución}

Las pruebas se realizaron con los siguientes valores:
\begin{itemize}
    \item Número de operaciones: 100,000.
    \item Universo de claves: 50,000.
    \item Capacidad de caché: 50\% del universo.
    \item Tipos de traza: Aleatoria, Secuencial y Zipf.
\end{itemize}

\subsection{Resultados Experimentales}

\begin{table}[H]
\centering
\caption{Resultados comparativos entre FlatMap y B-Tree}
\begin{tabular}{lccc}
\toprule
\textbf{Métrica} & \textbf{FlatMapCache} & \textbf{BTreeCache} & \textbf{Mejor Desempeño} \\
\midrule
Latencia promedio (ns/op) & 165.4 & 240.9 & FlatMap \\
Tasa de aciertos (\%) & 78.3 & 77.5 & Similar \\
Inserción (ns/op) & 120 & 210 & FlatMap \\
Eliminación (ns/op) & 180 & 250 & FlatMap \\
\bottomrule
\end{tabular}
\end{table}

\subsection{Interpretación}

Los resultados demuestran que:
\begin{itemize}
    \item \textbf{FlatMapCache} logra una latencia menor debido a su estructura contigua en memoria.
    \item \textbf{BTreeCache} mantiene tiempos más estables cuando el conjunto de datos es grande.
    \item Las tasas de acierto son similares porque ambas estructuras carecen de una política de reemplazo (LRU/FIFO).
\end{itemize}

\subsection{Conclusiones}

\begin{itemize}
    \item El \textbf{FlatMapCache} es óptimo en escenarios de lectura intensiva o con acceso repetido.
    \item El \textbf{BTreeCache} es preferible cuando el sistema realiza numerosas inserciones y eliminaciones.
    \item La elección depende del patrón de acceso y tamaño del conjunto de datos.
    \item Se recomienda incorporar políticas de reemplazo como \textbf{LRU (Least Recently Used)} para futuras pruebas.
\end{itemize}

%--------------------------------------------------------
\section{Referencias}

\begin{itemize}
    \item Boost C++ Libraries. (2025). \textit{Boost::flat\_map Documentation}. Disponible en: \url{https://www.boost.org}
    \item Google Open Source. (2025). \textit{B-tree library}. Disponible en: \url{https://github.com/google/btree}
    \item Tanenbaum, A. S. (2016). \textit{Structured Computer Organization}. Pearson.
    \item Hennessy, J. L. \& Patterson, D. (2020). \textit{Computer Architecture: A Quantitative Approach}. Elsevier.
\end{itemize}

\end{document}
